% Tradução completa para o português brasileiro da Introdução impressa (pp. físicas 3--4).
% A digitalização do IAS é a autoridade. NomiL e o inglês mantido servem para alinhamento.
% SGA-SRC-000052/000053/000055 são restaurações da impressão sem ramo em português.
% SGA-SRC-000054 é um erro genuíno de número no impresso e permanece visível por ramos.
% SGA-SRC-000059 é traduzido idiomaticamente nos dois modos; nenhum erro é inventado.

\chapter*{Introdução}

Este volume tem por objetivo facilitar ao não especialista o uso da
cohomologia $\ell$-ádica. Espero que muitas vezes lhe permita evitar o recurso
às densas exposições do SGA~4 e do SGA~5. Contém também alguns resultados
novos.

A primeira exposição, redigida por J.~F.~Boutot, oferece uma visão geral do
SGA~4. Apresenta os principais resultados ---com generalidade mínima, muitas
vezes insuficiente para as aplicações--- e uma ideia de sua demonstração.
Para obter resultados completos ou demonstrações detalhadas, o SGA~4 continua
sendo indispensável.

% SGA-SRC-000052: a impressão do IAS traz ``démonstration complète''; por isso,
% ambos os modos traduzem ``demonstração completa'', apesar do erro do NomiL.
O ``\nameref{II}'' contém uma demonstração completa da fórmula dos traços para
o endomorfismo de Frobenius. A demonstração é a apresentada por Grothendieck
no SGA~5, despojada de todo detalhe desnecessário.
% SGA-SRC-000053: a impressão do IAS restaura ``à l'utilisateur''; não há divisão em português.
Este Relatório deveria permitir ao usuário prescindir do SGA~5, que poderá ser
considerado como
\SGAreading{SGA-SRC-000054}{uma série de digressão, algumas muito interessantes}{uma série de digressões, algumas muito interessantes}.
A existência deste Relatório permitirá publicar em breve o SGA~5 tal como está.
Ele é completado pela exposição ``\nameref{VI}'', que explica como a fórmula
dos traços permite estudar somas trigonométricas e fornece exemplos.

O público ao qual se destinam as demais exposições é mais restrito, e seu
estilo reflete isso. A exposição ``\nameref{III}'' é uma generalização
``modular'' do Relatório, baseada no estudo das potências simétricas no SGA~4,
Exposição~XVII, \S~5.5.
% SGA-SRC-000055: a impressão do IAS traz o singular ``cette classe''; não há divisão.
A exposição ``\nameref{IV}'' define essa classe em diversos contextos e
estabelece a compatibilidade entre as interseções e os produtos cup. Em
``\nameref{V}'' reúnem-se alguns resultados conhecidos para os quais faltava
uma referência, bem como algumas compatibilidades. A exposição
``\nameref{VII}'' é nova. Em particular, ela fornece, para a cohomologia sem
suportes, teoremas de finitude análogos aos conhecidos para a cohomologia com
suportes compactos.

Para mais detalhes sobre as exposições, remeto às respectivas introduções.

% SGA-SRC-000059: o impresso ``étaien'' / NomiL ``étaien:'' é apenas evidência
% em português; ambos os modos apresentam o passado plural idiomático pretendido.
Por fim, agradeço a J.~L.~Verdier por ter me permitido reproduzir aqui suas
notas ``\nameref{VIII}''. Creio que continuam muito úteis e haviam se tornado
impossíveis de encontrar.

Nas referências internas deste volume, as exposições são citadas por um
título abreviado, indicado entre colchetes no sumário.

\vspace{20pt}
\begin{flushright}
Bures-sur-Yvette, 20 de setembro de 1976

\vspace{5pt}
Pierre Deligne
\end{flushright}
